% ============================================================ %
%  PREAMBLE                                                     %
%  Packages, fonts, spacing, and reusable macros only.          %
%  Per-document content lives in metadata.tex and template.tex. %
% ============================================================ %

% DOCUMENT CLASS SETUP %
\documentclass[12pt, a4paper]{article}

% DOCUMENT PACKAGES %
\usepackage[a4paper, margin=1in]{geometry}
\usepackage{enumitem}
\usepackage{amsmath}
\usepackage{amssymb}
\usepackage{mathtools}
\usepackage{graphicx}
\usepackage{xcolor}
\usepackage{listings}
\usepackage{hyperref}
\usepackage[capitalize, noabbrev]{cleveref}
\usepackage{csquotes}
\usepackage{tocloft}
\usepackage{array}
\usepackage{booktabs}
\usepackage{caption}
\usepackage{subcaption}
\usepackage{mwe}                 % provides \includegraphics{example-image*} placeholders
\usepackage{fancyhdr}
\usepackage{etoolbox}            % powers the optional bibliography toggle

% READABILITY / "AIRY" TYPOGRAPHY SETUP %
\usepackage{setspace}
\setstretch{1.15}
\setlength{\parindent}{0pt}
\setlength{\parskip}{0.6em plus 0.2em}
\usepackage{titlesec}
\titlespacing*{\section}{0pt}{2.4em}{1em}
\titlespacing*{\subsection}{0pt}{1.8em}{0.8em}
\titlespacing*{\subsubsection}{0pt}{1.4em}{0.6em}
\setlist{itemsep=0.4em, topsep=0.6em, parsep=0pt}

% FONT AND ENCODING SETUP %
\usepackage{fontspec}
\usepackage{polyglossia}
\usepackage[Thai]{ucharclasses}
\setdefaultlanguage{english}
\setotherlanguage{thai}
\setTransitionsFor{Thai}{\thaifont}{\normalfont}
\newfontfamily{\thaifont}{THSarabun.ttf}[
    Path=./fonts/,
    BoldFont=THSarabunBold.ttf,
    ItalicFont=THSarabunItalic.ttf,
    BoldItalicFont=THSarabunBoldItalic.ttf,
    Scale=1.23
]
\XeTeXlinebreaklocale "th"
\XeTeXlinebreakskip = 0pt plus 1pt

% BIBLIOGRAPHY (lightweight, optional per-document via \showbibliography) %
% Must load after polyglossia.
\usepackage[style=numeric, sorting=none]{biblatex}
\addbibresource{references.bib}

% CODE BLOCK PACKAGES %
\usepackage{tcolorbox}
\tcbuselibrary{minted, skins, breakable}

% DEFINE GITHUB DARK THEME COLORS %
\definecolor{ghback}{HTML}{0d1117}
\definecolor{ghtext}{HTML}{c9d1d9}
\definecolor{ghborder}{HTML}{30363d}

% CODE BLOCK CONFIGURATION %
\renewcommand{\theFancyVerbLine}{\textcolor{white!60!gray}{\scriptsize\arabic{FancyVerbLine}}}
\newtcblisting{codeblock}[3][]{
    enhanced,
    breakable,
    colback=ghback,
    colframe=ghborder,
    coltitle=ghtext,
    colupper=ghtext,
    fonttitle=\bfseries\small,
    arc=3mm,
    boxrule=0.5mm,
    % PADDING SETTINGS %
    left=25pt,
    right=10pt,
    top=5pt,
    bottom=5pt,
    % MINTED SETTINGS %
    listing engine=minted,
    listing only,
    minted language=#2,
    minted options={
        linenos,
        numbersep=5pt,
        fontsize=\footnotesize,
        breaklines=true,
        breakanywhere=true,
        autogobble,
        style=monokai,
    },
    title={#3},
    #1
}

% SEMANTIC CALLOUT BOXES %
% Kept to a small, muted set on purpose: Note, Definition, Example, Takeaway.
% Each takes an optional argument to override the default title text, e.g.
%   \begin{notebox}[Note: edge case] ... \end{notebox}
\newtcolorbox{notebox}[1][Note]{
    enhanced, breakable,
    colback=blue!4, colframe=blue!45!black,
    colbacktitle=blue!45!black, coltitle=white,
    fonttitle=\bfseries, title=#1,
    arc=1.5mm, boxrule=0.4mm,
    left=8pt, right=8pt, top=5pt, bottom=5pt,
}
\newtcolorbox{defbox}[1][Definition]{
    enhanced, breakable,
    colback=black!3, colframe=black!55,
    colbacktitle=black!55, coltitle=white,
    fonttitle=\bfseries, title=#1,
    arc=1.5mm, boxrule=0.4mm,
    left=8pt, right=8pt, top=5pt, bottom=5pt,
}
\newtcolorbox{examplebox}[1][Example]{
    enhanced, breakable,
    colback=green!4, colframe=green!30!black,
    colbacktitle=green!30!black, coltitle=white,
    fonttitle=\bfseries, title=#1,
    arc=1.5mm, boxrule=0.4mm,
    left=8pt, right=8pt, top=5pt, bottom=5pt,
}
\newtcolorbox{takeawaybox}[1][Key Takeaway]{
    enhanced, breakable,
    colback=orange!5, colframe=orange!55!black,
    colbacktitle=orange!55!black, coltitle=white,
    fonttitle=\bfseries, title=#1,
    arc=1.5mm, boxrule=0.4mm,
    left=8pt, right=8pt, top=5pt, bottom=5pt,
}

% DISABLE PAGE NUMBER IN TABLE OF CONTENTS %
\tocloftpagestyle{empty}

% RUNNING HEADER / FOOTER (body pages only; cover and TOC stay empty) %
% Left = course code, center = current section, right = author.
\pagestyle{fancy}
\setlength{\headheight}{14pt}
\fancyhf{}
\fancyhead[L]{\small\scshape \coursecode}
\fancyhead[C]{\small\itshape \rightmark}
\fancyhead[R]{\small \docauthor}
\fancyfoot[C]{\small\thepage}
\renewcommand{\headrulewidth}{0.4pt}
\renewcommand{\footrulewidth}{0pt}
\fancypagestyle{plain}{
    \fancyhf{}
    \fancyfoot[C]{\small\thepage}
    \renewcommand{\headrulewidth}{0pt}
}
